\documentclass{amsart}
% For dealing with references we use the comment environment
\usepackage{verbatim}
\newenvironment{reference}{\comment}{\endcomment}
%\newenvironment{reference}{}{}
\newenvironment{slogan}{\comment}{\endcomment}
\newenvironment{history}{\comment}{\endcomment}

% For commutative diagrams we use Xy-pic
\usepackage[all]{xy}

% We use 2cell for 2-commutative diagrams.
\xyoption{2cell}
\UseAllTwocells

% We use multicol for the list of chapters between chapters
\usepackage{multicol}

% This is generally recommended for better output
\usepackage{lmodern}
\usepackage[T1]{fontenc}

% For cross-file-references

% Package for hypertext links:
\usepackage{hyperref}

% For any local file, say "hello.tex" you want to link to please

% Theorem environments.
%
\theoremstyle{plain}
\newtheorem{theorem}[subsection]{Theorem}
\newtheorem{proposition}[subsection]{Proposition}
\newtheorem{lemma}[subsection]{Lemma}

\theoremstyle{definition}
\newtheorem{definition}[subsection]{Definition}
\newtheorem{example}[subsection]{Example}
\newtheorem{exercise}[subsection]{Exercise}
\newtheorem{situation}[subsection]{Situation}

\theoremstyle{remark}
\newtheorem{remark}[subsection]{Remark}
\newtheorem{remarks}[subsection]{Remarks}

\numberwithin{equation}{subsection}

% Macros
%
\def\lim{\mathop{\mathrm{lim}}\nolimits}
\def\colim{\mathop{\mathrm{colim}}\nolimits}
\def\Spec{\mathop{\mathrm{Spec}}}
\def\Hom{\mathop{\mathrm{Hom}}\nolimits}
\def\Ext{\mathop{\mathrm{Ext}}\nolimits}
\def\SheafHom{\mathop{\mathcal{H}\!\mathit{om}}\nolimits}
\def\SheafExt{\mathop{\mathcal{E}\!\mathit{xt}}\nolimits}
\def\Sch{\mathit{Sch}}
\def\Mor{\mathop{\mathrm{Mor}}\nolimits}
\def\Ob{\mathop{\mathrm{Ob}}\nolimits}
\def\Sh{\mathop{\mathit{Sh}}\nolimits}
\def\NL{\mathop{N\!L}\nolimits}
\def\CH{\mathop{\mathrm{CH}}\nolimits}
\def\proetale{{pro\text{-}\acute{e}tale}}
\def\etale{{\acute{e}tale}}
\def\QCoh{\mathit{QCoh}}
\def\Ker{\mathop{\mathrm{Ker}}}
\def\Im{\mathop{\mathrm{Im}}}
\def\Coker{\mathop{\mathrm{Coker}}}
\def\Coim{\mathop{\mathrm{Coim}}}

% Boxtimes
%
\DeclareMathSymbol{\boxtimes}{\mathbin}{AMSa}{"02}

%
% Macros for moduli stacks/spaces
%
\def\QCohstack{\mathcal{QC}\!\mathit{oh}}
\def\Cohstack{\mathcal{C}\!\mathit{oh}}
\def\Spacesstack{\mathcal{S}\!\mathit{paces}}
\def\Quotfunctor{\mathrm{Quot}}
\def\Hilbfunctor{\mathrm{Hilb}}
\def\Curvesstack{\mathcal{C}\!\mathit{urves}}
\def\Polarizedstack{\mathcal{P}\!\mathit{olarized}}
\def\Complexesstack{\mathcal{C}\!\mathit{omplexes}}
% \Pic is the operator that assigns to X its picard group, usage \Pic(X)
% \Picardstack_{X/B} denotes the Picard stack of X over B
% \Picardfunctor_{X/B} denotes the Picard functor of X over B
\def\Pic{\mathop{\mathrm{Pic}}\nolimits}
\def\Picardstack{\mathcal{P}\!\mathit{ic}}
\def\Picardfunctor{\mathrm{Pic}}
\def\Deformationcategory{\mathcal{D}\!\mathit{ef}}

\setlength{\emergencystretch}{3em}
\begin{document}
\title[Stacks Project, Tag 017K]{705 --- Finite-limit representability for finite eventually degenerate simplicial sets}
\maketitle
\noindent Copyright (C) 2005--2025 Johan de Jong. Stacks Project authors. Source: \href{https://stacks.math.columbia.edu/tag/017K}{Tag 017K}.

\noindent Editorial difficulty note (not author text). Difficulty H2: The proof constructs a finite initial subcategory of the simplex indexing category using simultaneous degeneracy removal and a dimension bound. Establishing finite representability from this combinatorial cofinality argument is above ordinary qualification category theory..

\noindent Editorial provenance note (not author text). History: excerpt from revision \texttt{a0444\allowbreak 6e57e\allowbreak c1fbc\allowbreak 252a8\allowbreak 71afc\allowbreak ec775\allowbreak 2fb28\allowbreak 07b14}, simplicial.tex, lines 1460--1543. Reference presentation was adapted; original-author context and core support blocks were retained or added. The mathematical wording is unchanged. Licensed under GNU FDL 1.2 or later; no invariant sections or cover texts. See ../licenses/COPYING. Context blocks reproduce author definitions and supporting results; they are not additional counted problems.

\begin{definition}
\label{categories-definition-initial}
Let $H : \mathcal{I} \to \mathcal{J}$ be a functor between categories.
We say {\it $\mathcal{I}$ is initial in $\mathcal{J}$} or that
$H$ is {\it initial} if
\begin{enumerate}
\item for all $y \in \Ob(\mathcal{J})$ there exist an
$x \in \Ob(\mathcal{I})$ and a morphism $H(x) \to y$,
\item for any $y \in \Ob(\mathcal{J})$, $x , x' \in \Ob(\mathcal{I})$ and
morphisms $H(x) \to y$, $H(x') \to y$ in $\mathcal{J}$
there exist a sequence of morphisms
$$
x = x_0 \leftarrow x_1 \rightarrow x_2 \leftarrow x_3 \rightarrow \ldots
\rightarrow x_{2n} = x'
$$
in $\mathcal{I}$ and morphisms $H(x_i) \to y$ in $\mathcal{J}$
such that the diagrams
$$
\xymatrix{
H(x_{2k}) \ar[rd] &
H(x_{2k + 1}) \ar[l] \ar[r] \ar[d] &
H(x_{2k + 2}) \ar[ld] \\
& y
}
$$
commute for $k = 0, \ldots, n - 1$.
\end{enumerate}
\end{definition}

\begin{lemma}
\label{lemma-exists-hom-0-from-simplicial-set}
Assume the category $\mathcal{C}$
has coproducts of any two objects and countable
limits. Let $U$ be a simplicial set, with $U_n$ finite nonempty
for all $n \geq 0$.
Let $V$ be a simplicial object of $\mathcal{C}$.
Then the functor
\begin{eqnarray*}
\mathcal{C}^{opp} & \longrightarrow & \textit{Sets} \\
X
& \longmapsto &
\Mor_{\text{Simp}(\mathcal{C})}(X \times U, V)
\end{eqnarray*}
is representable.
\end{lemma}

\begin{proof}
A morphism from $X \times U$ into $V$ is given by a collection
of morphisms $f_u : X \to V_n$ with $n \geq 0$ and $u \in U_n$.
And such a collection actually defines a morphism if and only
if for all $\varphi : [m] \to [n]$ all the diagrams
$$
\xymatrix{
X \ar[r]^{f_u} \ar[d]_{\text{id}_X} & V_n \ar[d]^{V(\varphi)} \\
X \ar[r]^{f_{U(\varphi)(u)}} & V_m
}
$$
commute. Thus it is natural to introduce a category
$\mathcal{U}$ and a functor
$\mathcal{V} : \mathcal{U}^{opp} \to \mathcal{C}$
as follows:
\begin{enumerate}
\item The set of objects of $\mathcal{U}$ is
$\coprod_{n \geq 0} U_n$,
\item a morphism from $u' \in U_m$ to $u \in U_n$
is a $\varphi : [m] \to [n]$ such that $U(\varphi)(u) = u'$
\item for $u \in U_n$ we set $\mathcal{V}(u) = V_n$, and
\item for $\varphi : [m] \to [n]$ such that $U(\varphi)(u) = u'$
we set $\mathcal{V}(\varphi) = V(\varphi) : V_n \to V_m$.
\end{enumerate}
At this point it is clear that our functor is nothing but the
functor defining
$$
\lim_{\mathcal{U}^{opp}} \mathcal{V}
$$
Thus if $\mathcal{C}$ has countable limits then this limit
and hence an object representing the functor of the lemma
exist.
\end{proof}


\begin{lemma}
\label{lemma-exists-hom-0-from-simplicial-set-finite}
Assume the category $\mathcal{C}$
has coproducts of any two objects and finite
limits. Let $U$ be a simplicial set, with $U_n$ finite nonempty
for all $n \geq 0$. Assume that all $n$-simplices
of $U$ are degenerate for all $n \gg 0$.
Let $V$ be a simplicial object of $\mathcal{C}$.
Then the functor
\begin{eqnarray*}
\mathcal{C}^{opp} & \longrightarrow & \textit{Sets} \\
X
& \longmapsto &
\Mor_{\text{Simp}(\mathcal{C})}(X \times U, V)
\end{eqnarray*}
is representable.
\end{lemma}

\begin{proof}
We have to show that the category $\mathcal{U}$ described
in the proof of Lemma \ref{lemma-exists-hom-0-from-simplicial-set}
has a finite subcategory $\mathcal{U}'$ such that the limit
of $\mathcal{V}$ over $\mathcal{U}'$ is the same as the
limit of $\mathcal{V}$ over $\mathcal{U}$. We will use
Categories, Lemma \href{https://stacks.math.columbia.edu/tag/002R}{lemma initial (Tag 002R)}.
For $m > 0$ let $\mathcal{U}_{\leq m}$ denote the full
subcategory with objects $\coprod_{0 \leq n \leq m} U_m$.
Let $m_0$ be an integer such that every $n$-simplex
of the simplicial set $U$ is degenerate if $n > m_0$.
For any $m \geq m_0$ large enough, the subcategory
$\mathcal{U}_{\leq m}$ satisfies property (1) of
Categories, Definition \ref{categories-definition-initial}.

\medskip\noindent
Suppose that $u \in U_n$ and
$u' \in U_{n'}$ with $n, n' \leq m_0$ and suppose that
$\varphi : [k] \to [n]$, $\varphi' : [k] \to [n']$
are morphisms such that $U(\varphi)(u) = U(\varphi')(u')$.
A simple combinatorial argument shows that if $k > 2m_0$,
then there exists an index $0 \leq i \leq 2m_0$ such that
$\varphi(i) =\varphi(i + 1)$ and $\varphi'(i) = \varphi'(i + 1)$.
(The pigeon hole principle would tell you this works if
$k > m_0^2$ which is good enough for the argument below
anyways.) Hence, if $k > 2m_0$, we may write
$\varphi = \psi \circ \sigma^{k - 1}_i$ and
$\varphi' = \psi' \circ \sigma^{k - 1}_i$ for some
$\psi : [k - 1] \to [n]$ and some $\psi' : [k - 1] \to [n']$.
Since $s^{k - 1}_i : U_{k - 1} \to U_k$ is injective,
see Lemma \href{https://stacks.math.columbia.edu/tag/016F}{lemma si injective (Tag 016F)}, we conclude that
$U(\psi)(u) = U(\psi')(u')$ also. Continuing in this
fashion we conclude that given morphisms
$u \to z$ and $u' \to z$ of $\mathcal{U}$
with $u, u' \in \mathcal{U}_{\leq m_0}$, there exists
a commutative diagram
$$
\xymatrix{
u \ar[rd] \ar[rrd] & & \\
& a \ar[r] & z \\
u' \ar[ru] \ar[rru]
}
$$
with $a \in \mathcal{U}_{\leq 2m_0}$.

\medskip\noindent
It is easy to deduce from this that the finite subcategory
$\mathcal{U}_{\leq 2m_0}$ works. Namely, suppose given
$x' \in U_n$ and $x'' \in U_{n'}$ with $n, n' \leq 2m_0$ as well as
morphisms $x' \to x$ and $x'' \to x$ of $\mathcal{U}$
with the same target. By our choice of $m_0$ we can
find objects $u, u'$ of $\mathcal{U}_{\leq m_0}$ and
morphisms $u \to x'$, $u' \to x''$.
By the above we can find $a \in \mathcal{U}_{\leq 2m_0}$
and morphisms $u \to a$, $u' \to a$ such that
$$
\xymatrix{
u \ar[rd] \ar[rrd] \ar[r] & x' \ar[rd] & \\
& a \ar[r] & x \\
u' \ar[ru] \ar[rru] \ar[r] & x'' \ar[ru] &
}
$$
is commutative. Turning this diagram 90 degrees clockwise
we get the desired diagram as in (2) of
Categories, Definition \ref{categories-definition-initial}.
\end{proof}
\end{document}
